\documentclass[11pt]{article}
\usepackage[a4paper,margin=27mm]{geometry}
\usepackage{amsmath,amssymb,mathtools}
\usepackage{fontspec,unicode-math}
\setmainfont{Latin Modern Roman}
\setmathfont{Latin Modern Math}
\usepackage{xurl}
\usepackage[hidelinks]{hyperref}
\newcommand{\sref}[2]{\hyperlink{source:#1:#2}{[S#2]}}
\newcommand{\cataloguescope}[1]{\paragraph{Mathematical question.}#1}
\hypersetup{pdftitle={Problem statement}}
\begin{document}
\section{Complex continuation of the associated zeta function}\label{statement}

\subsection{Definitions and mathematical statement}

\paragraph{Shared notation.}$\mathbb N=\{1,2,\ldots\}$, and $\mathbb Z,\mathbb Q,\mathbb R,\mathbb C$ denote the integers, rationals, reals and complex numbers. Any other conventions are specified locally.


Let
\[
\mathcal A=\{(a,b,c,d)\in\mathbb Z_{\ge0}^{4}:ad-bc=1\},\qquad
f(a,b,c,d)=\sqrt{a^2+b^2}+\sqrt{c^2+d^2}-\sqrt{(a+c)^2+(b+d)^2}.
\]
Every such $f$ is positive, so its complex power means $f^s=\exp(s\log f)$ with the real logarithm. Define
\[
Z(s)=\sum_{(a,b,c,d)\in\mathcal A} f(a,b,c,d)^s.
\]
The original published question cites a convergence baseline $s>1/2$. The authors' later arXiv version corrects this and records convergence for real $s>2/3$ and divergence for $s\le2/3$. The same positive-term estimate gives absolute and locally uniform convergence for $\operatorname{Re}s>2/3$, which is the initial holomorphic domain used here. \sref{statement}{3}
\paragraph{Statement recorded in the supplied source.}
Extend the function $Z(s)$ to complex values of $s$: determine an analytic continuation beyond its initial half-plane of convergence and the complex domains on which such a continuation exists. \sref{statement}{2}
\subsection{Short English statement}
Continue the lattice sum of complex powers of triangle-inequality defects beyond its convergent half-plane.
\subsection{Sources}
\begin{description}
\item[\hypertarget{source:statement:1}{S1}] ULAM.AI and the UnsolvedMath Contributors, UnsolvedMath: A Curated Collection of Open Mathematics Problems, catalogue entry. CC BY 4.0. \url{https://huggingface.co/datasets/ulamai/UnsolvedMath}.
\item[\hypertarget{source:statement:2}{S2}] Nikita Kalinin and Mikhail Shkolnikov, \emph{The number $\pi$ and summation by $SL(2,\mathbb Z)$}, Arnold Mathematical Journal (2018) \S4, Questions, pp. 5--6. \url{https://amj.math.stonybrook.edu/html-articles/Files-2015-2024/17-75/index.html}.
\item[\hypertarget{source:statement:3}{S3}] Nikita Kalinin and Mikhail Shkolnikov, \emph{Tropical formulae for summation over a part of $SL(2,\mathbb Z)$}, arXiv:1711.02089 (updated version) \S1.6, Petrov's comment and Remark 7, pp. 6--7. \url{https://arxiv.org/pdf/1711.02089}.
\end{description}
\label{end:statement}
\end{document}
